% Teil II: der durchgerechnete Weg.

\subsection*{Gegeben und gesucht}
Aussage: ``Zu jeder beliebigen Zahl gibt es eine natürliche Zahl, die
größer ist.'' Gesucht: die Negation aus den vier vorgegebenen Aussagen.

\subsection*{Ansatz}
Negieren der Quantoren ``zu jeder'' $\rightarrow$ ``es gibt eine'', ``es
gibt eine ($> x$)'' $\rightarrow$ ``keine ($> x$)''.

\subsection*{Rechnung}
Formal: $\bigl(\forall x\,\exists n\in\mathbb N:\ n>x\bigr)$. Die Negation:
$\exists x$ so, dass $\forall n\in\mathbb N:\ n\not>x$, d.h. $n\leq x$.
In Worten: \emph{Es gibt mindestens eine Zahl, die von keiner natürlichen
Zahl übertroffen wird.} Das trifft auf Aussage b) zu.

\begin{ergebnis}
Die Negation ist \textbf{b)}: ``Es gibt mindestens eine Zahl, die von
keiner natürlichen Zahl übertroffen wird.''
\end{ergebnis}

\begin{kontrolle}
Aussage a) ``keine einzige Zahl wird von jeder übertroffen'' ist zu stark
(verneint nur den Allquantor, nicht den Existenzquantor); c) ersetzt
``größer'' durch ``kleiner''; d) verneint den Existenzteil falsch. Nur b)
entspricht der formalen Negation.
\end{kontrolle}

\begin{merksatz}
Die Negation dreht jeden Quantor um (``jede'' $\leftrightarrow$
``mindestens eine'', ``es gibt'' $\leftrightarrow$ ``keine''); der
Wahrheitsgehalt spielt keine Rolle.
\end{merksatz}
